% Part: computability
% Chapter: computability-theory
% Section: computable-sets

\documentclass[../../../include/open-logic-section]{subfiles}

\begin{document}

\olfileid{cmp}{thy}{cps}
\olsection{संगणनक्षम संच}

संगणनक्षमतेची कल्पना संगणनक्षम फलनांपासून संगणनक्षम संचांपर्यंत
विस्तारता येते:

\begin{defn}
  $S$ हा नैसर्गिक संख्यांचा संच असू द्या. त्याचे जातिबोधक
  फलन~$\Char{S}$ संगणनक्षम असेल तेव्हाच $S$ \emph{संगणनक्षम}
  असतो. दुसऱ्या शब्दांत, पुढील फलन संगणनक्षम असेल तेव्हाच
  $S$~संगणनक्षम असतो:
\[
\Char{S}(x) =
\begin{cases}
1 & \text{जर $x \in S$ असेल तर} \\
0 & \text{अन्यथा}
\end{cases}
\]
त्याचप्रमाणे, संबंध $R(x_0, \dots, x_{k-1})$ संगणनक्षम असतो
तेव्हाच त्याचे जातिबोधक फलन संगणनक्षम असते.

संगणनक्षम संच आणि संबंध यांना \emph{निर्णेय} असेही म्हणतात.
\end{defn}

\begin{explain}
आता संगणनक्षमतेच्या अनेक कल्पना आहेत: आंशिक फलनांसाठी, फलनांसाठी
आणि संचांसाठी. त्यांची गल्लत करू नका!{} \iftag{TMs}{आंशिक
  फलनाचे संगणन करणारे ट्यूरिंग यंत्र, ज्या आदानांवर फलन परिभाषित
  आहे त्यांसाठी फलनाचे निर्गत परत करते; संचाचे संगणन करणारे ट्यूरिंग
  यंत्र आदान संचात आहे की नाही हे ठरवून $1$ किंवा~$0$ परत करते.}{}
\end{explain}

\end{document}
